% Part: sets-functions-relations
% Chapter: sets
% Section: russells-paradox

\documentclass[../../../include/open-logic-section]{subfiles}


\begin{document}

\olfileid{sfr}{set}{rus}
\olsection{Paradoks Russell}

Ekstensionalitas mbenerake notasi $\Setabs{x}{\phi(x)}$ kanggo
\emph{himpunan kasebut} kang ngemot kabeh $x$ kang ndadekake~$\phi(x)$.
Nanging, kang \emph{tenan} dibenerake dening ekstensionalitas
mung gagasan iki: \emph{yen} ana himpunan kang anggotane kabeh
lan mung objek kang nduweni sipat $\phi$, \emph{mula} himpunan
kaya mangkono mung ana siji. Kanthi tembung liya: sawise netepake
sawijining~$\phi$, himpunan $\Setabs{x}{\phi(x)}$ iku tunggal,
\emph{yen ana}.

Nanging pratelan mawa syarat iki wigati!{} Ora saben sipat
bisa digunakake kanggo \emph{komprehensi}. Tegese, ana sipat
kang \emph{ora} nemtokake himpunan. Yen kabeh sipat nemtokake
himpunan, kita bakal nemoni kontradiksi kang cetha.
Tuladha kang paling misuwur yaiku Paradoks Russell.

Himpunan bisa dadi !!{element}s saka himpunan liyane---tuladhane,
himpunan pangkat saka himpunan~$A$ dumadi saka himpunan-himpunan.
Mula ana tegese yen kita takon utawa nliti apa sawijining
himpunan iku !!a{element} saka himpunan liyane.
Apa himpunan bisa dadi anggota awake dhewe? Katone ora ana
apa-apa ing gagasan himpunan kang ngalangi iki.
Tuladhane, yen \emph{kabeh} himpunan mbentuk kumpulan objek,
kita bisa ngira yen kabeh mau bisa diklumpukake dadi siji
himpunan---himpunan kabeh himpunan. Lan himpunan kasebut,
amarga dhewe uga himpunan, bakal dadi !!a{element}
saka himpunan kabeh himpunan.

Paradoks Russell tuwuh nalika kita ngrembug sipat ora duwe
awake dhewe minangka !!a{element}, yaiku \emph{ora dadi anggota
awake dhewe}. Kepriye yen kita nganggep ana himpunan kabeh
himpunan kang ora duwe awake dhewe minangka !!a{element}?
Apa
\[
R = \Setabs{x}{x \notin x}
\]
ana? Pranyata kita bisa mbuktekake yen himpunan iki ora ana.

\begin{thm}[Paradoks Russell]\ollabel{thm:russells-paradox}
	Ora ana himpunan $R = \Setabs{x}{x \notin x}$.
\end{thm}

\begin{proof}
Yen $R = \Setabs{x}{x \notin x}$ ana, mula
$R \in R$ yen lan mung yen $R \notin R$, kang arupa kontradiksi.
\end{proof}

\begin{tagblock}{novice}
\begin{explain}
Ayo nlusuri bukti iki kanthi luwih alon. Yen $R$ ana,
pitakon apa $R \in R$ utawa ora iku ana tegese.
Umpamakna pancen $R \in R$. $R$~wis ditetepake minangka
himpunan kabeh himpunan kang ora dadi !!{element}s saka
awake dhewe. Mula, yen $R \in R$, $R$ dhewe ora duwe
sipat kang digunakake kanggo nemtokake $R$.
Nanging mung himpunan kang duwe sipat iki kang ana ing~$R$,
mula $R$ ora bisa dadi !!a{element} saka~$R$,
yaiku $R \notin R$. Nanging $R$ ora bisa bebarengan dadi
lan ora dadi !!a{element} saka~$R$, mula kita oleh kontradiksi.

Amarga anggepan yen $R \in R$ nuwuhake kontradiksi,
kita oleh $R \notin R$. Nanging iki uga nuwuhake kontradiksi!{}
Amarga yen $R\notin R$, $R$ dhewe pancen duwe sipat kang
digunakake kanggo nemtokake $R$, mula $R$ bakal dadi
!!a{element} saka $R$, kaya kabeh himpunan liyane kang
ora dadi anggota awake dhewe.
Lan maneh, ora bisa bebarengan ora dadi lan dadi
!!a{element} saka~$R$.
\end{explain}
\end{tagblock}

\begin{digress}
Kepriye carane mbangun teori himpunan kang ora kena Paradoks
Russell, yaiku kang ora nggawe pratelan \emph{ora konsisten}
yen $R = \Setabs{x}{x \notin x}$ ana?
Kita kudu netepake aksioma kang menehi syarat kanthi tliti banget
kanggo nyatakake kapan himpunan ana (lan kapan ora ana).

Teori himpunan kang digarisake ing bab iki ora nindakake iku.
Teori iki \emph{pancen naif}. Teori iki mung nerangake yen
himpunan manut ekstensionalitas lan yen kita duwe sawetara
himpunan, kita bisa mbentuk gabungan, irisan, lan sapanunggalane.
Teori himpunan bisa dikembangake kanthi luwih tliti miturut
paugeran matematis tinimbang iki.
\oliflabeldef{cumul:::part}{Andharan kanthi tliti kasebut diwenehake
ing Perangan \olref[cumul][][]{part}. Saiki kita bakal nerusake
kanthi cara naif, lan kanthi ngati-ati ngupaya nyingkiri kontradiksi.}{}
\end{digress}


\end{document}
